% !TEX root = ../main.tex
\clearpage
\section{Bitácora de decisiones}\label{sec:bitacora}

La Tabla~\ref{tab:bitacora} registra las decisiones relevantes del proyecto, cada una con la alternativa descartada y la razón que la sustenta.

{\tablacompacta
\begin{longtable}{@{}L{2.05cm}L{3.8cm}L{3.3cm}L{5.15cm}@{}}
  \caption{Bitácora de decisiones de diseño}\label{tab:bitacora}\\
  \toprule
  \cab{Fecha} & \cab{Decisión} & \cab{Alternativa descartada} & \cab{Razón} \\ \midrule \endfirsthead
  \toprule
  \cab{Fecha} & \cab{Decisión} & \cab{Alternativa descartada} & \cab{Razón} \\ \midrule \endhead
  \bottomrule \endlastfoot
  18/09/2026 & Dominó como juego del proyecto, desarrollado en \Lenguaje. & Otros juegos de mesa. & Ofrece estados, variantes y jugadas reversibles; \Lenguaje\ es el lenguaje de mayor dominio de mi parte. \\ \filasep
  27/09/2026 & Partida local de un jugador humano contra una a tres CPU. & Multijugador en red con salas y código de acceso. & La red concentra el esfuerzo en sincronización, desconexiones e información oculta, aspectos que no se evalúan; además consume el presupuesto de clases y dificulta el modo banco de pruebas. La abstracción del jugador conserva la posibilidad de agregar un jugador remoto sin modificar el dominio. \\ \filasep
  27/09/2026 & Juego siempre individual. & Modalidad en parejas. & Las parejas solo aplican con cuatro jugadores y agregan reglas de puntuación compartida sin aportar fuerzas de diseño nuevas. \\ \filasep
  27/09/2026 & Tipos de dominó doble~6 y doble~9. & Doble~12. & No aporta un comportamiento distinto al de doble~9; únicamente prolonga la partida. \\ \filasep
  27/09/2026 & Abre la mano quien tenga el doble más alto entre los jugadores (RN-05). & Exigir siempre el doble máximo del juego. & Cuando quedan fichas para robar, el doble máximo puede quedar sin repartir y la mano no podría comenzar. \\ \filasep
  27/09/2026 & Deshacer solo la última jugada del humano, junto con las jugadas posteriores de las CPU. & Deshacer varios niveles y rehacer. & Permite revertir jugadas con menor complejidad; el historial permitiría extenderlo más adelante. \\ \filasep
  27/09/2026 & Partida a una sola mano, sin puntos acumulados. & Varias manos con meta de puntos. & Simplifica el cierre de la partida; los puntos se utilizan únicamente para resolver el juego cerrado. \\ \filasep
  27/09/2026 & Límite de tiempo opcional por turno (10, 20 o 30 s); al agotarse, se pierde el turno. & Que el sistema juegue o robe por el jugador. & Es la regla más simple de comunicar y no decide en nombre del jugador. \\ \filasep
  27/09/2026 & Incluir guardar y cargar partida. & No persistir la partida. & Completa el ciclo de vida de la jugada registrada y permite medir operaciones de disco. \\ \filasep
  27/09/2026 & Acoger el componente de medición. & No presentarlo. & Permite sustentar con datos el costo de las decisiones de diseño. \\ \filasep
  02/10/2026 & Guardar la partida como configuración, semilla y jugadas, y reconstruirla al cargar. & Serializar el estado completo de \texttt{Partida}. & No obliga a exponer el estado interno y reutiliza el historial que ya exige deshacer; el costo es repetir las jugadas al cargar, que se medirá como acción crítica. \\ \filasep
  02/10/2026 & \texttt{Jugada} identifica al jugador por su posición. & Referenciar el objeto \texttt{Jugador}. & La jugada es un valor inmutable que se guarda en archivo; la posición la identifica sin acoplarla a la entidad. \\ \filasep
  02/10/2026 & Tipos de valor inmutables (\texttt{record}) para fichas, jugadas, resultados, configuración y vista. & Clases mutables relacionadas por asociación. & Se comparan por valor, pueden compartirse sin copias defensivas y la vista no permite modificar el dominio. \\ \filasep
  02/10/2026 & \texttt{Partida} recibe \texttt{IReloj} y aplica la regla de tiempo. & Calcular el tiempo en \texttt{ControladorJuego}. & RN-04 es una regla del juego; ubicarla en el controlador le restaría cohesión. \\ \filasep
  02/10/2026 & Proyecto \texttt{Domino.Persistencia} para el repositorio JSON. & Repositorio dentro de \texttt{Domino.Consola}. & La simulación también guarda y carga partidas para medir esas acciones. \\ \filasep
  02/10/2026 & Un solo repositorio con el código y el documento. & Repositorios separados. & Una misma etiqueta marca el estado del código y del diseño de cada entrega, y el historial muestra los cambios de ambos juntos. \\ \filasep
  02/10/2026 & En la simulación, el simulador ocupa el lugar del humano. & Partidas solo entre CPU. & La regla RN-02 exige un jugador humano; simularlo conserva la regla y ejercita el mismo código que la consola. \\ \filasep
  24/10/2026 & Strategy para el criterio de decisión de cada nivel de CPU. & Mantener las subclases de \texttt{JugadorCpu}; Template Method. & Los niveles solo difieren en la elección completa; la estrategia se agrega sin tocar el jugador ni la partida (decisión D1). \\ \filasep
  24/10/2026 & State para las reglas que dependen del estado de la partida. & Condicionales sobre \texttt{EstadoPartida}; Chain of Responsibility. & Agrupa las reglas de cada estado y hace explícitas las transiciones (decisión D2). \\ \filasep
  24/10/2026 & Command: \texttt{Jugada} pasa de tipo de valor a jerarquía de comandos. & Memento; Prototype. & La misma secuencia de jugadas sirve para deshacer, mostrar y guardar sin exponer el estado interno (decisión D3). \\ \filasep
  24/10/2026 & Retirar el estado Cerrada. & Modelarlo como estado transitorio antes de Terminada. & Duraba un instante y no tenía reglas propias; la forma de cierre ya queda en \texttt{Resultado}. \\ \filasep
  24/10/2026 & El límite de tiempo no corre durante la salida (RN-04). & Aplicarlo también a la salida. & Si quien abre pierde el turno, nadie más puede abrir; además, la ficha de salida es una jugada obligada. \\ \filasep
  24/10/2026 & \texttt{IniciarPartida} y \texttt{CargarPartida} resuelven los turnos de las CPU. & Esperar la primera acción del humano. & Si la salida correspondía a una CPU, la partida quedaba detenida sin opciones para el humano. \\ \filasep
  24/10/2026 & No aplicar Chain of Responsibility a la validación ni Observer a la notificación de jugadas. & Aplicarlos en esta entrega. & No hay una fuerza que los justifique: validaciones fijas y un único interesado en las jugadas (decisiones D4 y D5). \\ \filasep
  24/10/2026 & Registrar los tipos de jugada para el JSON en \texttt{Repositorio\allowbreak Partidas\allowbreak Json}. & Atributos de serialización en las clases del dominio. & El dominio no depende del formato del archivo. \\ \filasep
  24/10/2026 & Consola interactiva: un solo tablero redibujado, fichas con sus puntos y manejo con teclado, con \texttt{PantallaConsola} separada de \texttt{InterfazConsola}. & Consola por líneas que se desplaza. & Mejora la jugabilidad sin tocar el dominio; separar el dibujo de la lectura de teclas conserva la cohesión de cada clase. \\ \filasep
  24/10/2026 & El turno se pierde en el instante en que se agota el tiempo: la consola consulta \texttt{VerificarTiempo} mientras espera una tecla. & Verificar el tiempo solo al recibir la acción. & Cumple RN-04 sin que el jugador tenga que presionar una tecla para descubrir que perdió el turno. \\ \filasep
  24/10/2026 & Proyecto de pruebas automatizadas \texttt{Domino.Pruebas}. & Verificar solo jugando en la consola. & Verifica las reglas y los patrones en cada cambio. No forma parte del diseño ni del conteo de clases. \\ \filasep
  30/10/2026 & Presentar la línea base de medición en la Entrega~3. & Presentarla en esta entrega. & Se priorizaron los patrones obligatorios; la simulación ya ejercita las acciones críticas sobre el escenario que se medirá. \\
\end{longtable}}

\subsection{Declaración de uso de asistentes de inteligencia artificial}

En la Entrega~1 se utilizó un asistente de inteligencia artificial para explorar alternativas de alcance del juego, revisar la redacción del caso de uso CU-01 frente a la plantilla del curso, proponer borradores del modelo de dominio y del diagrama de clases, mejorar la redacción de las secciones de este documento, organizar su estructura en LaTeX y revisar el diagrama de clases frente a las reglas de negocio. Las decisiones registradas en esta bitácora, incluidas las alternativas descartadas, fueron de análisis propio.

En la Entrega~2 se utilizó un asistente de inteligencia artificial para analizar la pertinencia de los patrones escogidos frente al juego y como apoyo en la escritura del documento.
